% Proposed replacement section. This is a written metatheorem, not a Lean proof.
% Baseline: ProveIt d9383325b18d349b8bfdc54ef564c5784cc029ec.
\section{Introduction selected by a reflexive base}
\label{sec:reflexive-base}

The native zero-normal calculus admits a single introduction whose
selection test uses the formula it would otherwise retain as the base.
It requires no distinguished witness and no syntactically disguised zero.
The construction is an instance of the guarded-introduction method, with
a useful difference: its guard need only be preserved by substitution;
it need not be reflected.  A newly reflexive base permits an elementary
repair of the substituted inference.

\subsection{Normal syntax and the complete rule inventory}

Let $\operatorname{nf}$ be the existing structural zero normalization:
it recursively normalizes the children of a term, then replaces
$I(0,a;u.h)$ by $a$.  It performs no successor computation.
Terms and formulas in this section are zero-normal.  Every displayed
substitution means capture-free substitution followed by
$\operatorname{nf}$ on all numerical terms.  Write $B(t,r)$ for this
normal instance of a matrix whose distinguished coordinates are $x,y$.
The formula constructors are still equality, $Q$, and $D$; their ordinary
readings are unchanged:
\[
 Q_yB=B(0)\to\forall y\,B(y),\qquad
 D_{x,y}B=\exists x\bigl(B(x,0)\to\forall y\,B(x,y)\bigr).
\]
The numerical syntax is $x,0,S(t),I(n,a;u.h)$, with the literal-zero
counter excluded from a normal iterator node.  The equation $0=S0$ is
the designated contradiction.

Define the decidable metasyntactic predicate
\[
 \mathsf R(B,t)\quad\Longleftrightarrow\quad
 B(t,0)\text{ is literally }r=r\text{ for some normal term }r.
\]
This is a comparison of two finite syntax trees, not a call to arithmetic
equality or a theorem prover.  Since substitution preserves the formula
root, $\mathsf R(B,t)$ entails that $B$ is a literal equation.
Set
\[
 \eta_{B,t}(z)=
 \begin{cases}z,&\mathsf R(B,t),\\0,&\text{otherwise}.\end{cases}
\]
The calculus $\mathcal R$ has exactly these six primitive rules:
\begin{align*}
 &\frac{A\in\Gamma}{\Gamma\vdash A} &&\mathrm{Hyp},\\[3pt]
 &\Gamma\vdash
 I(Sn,a;u.h)=h\bigl(I(n,a;u.h)\bigr) &&\mathrm{IS},\\[3pt]
 &\frac{\Gamma\vdash s=t\qquad\Gamma\vdash h(s)=r}
        {\Gamma\vdash r=h(t)} &&\mathrm{Tr},\\[5pt]
 &\frac{\Gamma\vdash Q_yB\qquad\Gamma\vdash B(0)}
        {\Gamma\vdash B(t)} &&\mathrm{QE},\\[5pt]
 &\frac{\Gamma\vdash D_{x,y}B\qquad
         \Gamma,Q_yB(b,y)\vdash A}{\Gamma\vdash A} &&\mathrm{DE},\\[5pt]
 &\frac{\Gamma,B(t,\eta_{B,t}(z))\vdash B(t,Sz)}
        {\Gamma\vdash D_{x,y}B} &&\mathrm{DI}_{\mathsf R}.
\end{align*}
All terms in IS and Tr are understood in the normal syntax, including
the possible contraction of $I(n,a;u.h)$ when $n$ is literally zero.
In DE, $b$ is fresh for the retained context, conclusion, and fixed
parameters.  In DI, $z$ is fresh for the retained context, witness, and
fixed parameters.  No other equality, logical, structural, induction,
or computation rule is primitive.

For comparison let $\mathcal Z$ be the repository's seven-rule
zero-normal calculus: it has the first five rules above and, instead of
$\mathrm{DI}_{\mathsf R}$, has unrestricted D introduction by cases and
Q introduction by equation progress:
\[
 \frac{\Gamma,B(t,0)\vdash B(t,Sz)}{\Gamma\vdash D_{x,y}B},
 \qquad
 \frac{\Gamma,E(z)\vdash E(Sz)}{\Gamma\vdash Q_yE(y)}.
\]
In the second rule $E$ is a literal equation.  The repository proves
that $\mathcal Z$ derives the normal form of exactly the native sequents.

\subsection{The noncircular structural bootstrap}

\paragraph{Normalization facts used.}
The existing normalization layer proves idempotence and compatibility
with substitution.  In particular the normal form of a substitution
instance depends only on the normal forms of the term and the
substitution's values.  Thus normal substitution is associative with
the usual lifted substitutions under binders.  These facts are
syntactic theorems; they do not use any rule of $\mathcal R$.

\paragraph{Weakening, renaming, and cut come first.}
Weakening, exchange, contraction, injective renaming of free variables,
and simultaneous ordinary-assumption cut are admissible by induction on
derivations.  These operations leave the matrix and witness unchanged,
or apply an injective variable renaming.  A reflexive equation stays
reflexive and a nonreflexive equation cannot become reflexive under
such a renaming.  Hence no DI mode changes in these inductions.
At rules with eigenvariables first choose fresh names away from the
inserted derivations.  This argument does not yet invoke arbitrary
numerical substitution.

\paragraph{Equality comes before missing cases.}
For any normal term $r$, let
\[
 J_r:=I(S0,r;u.u).
\]
IS gives $J_r=r$: its right-hand side is the zero-normal form of
$I(0,r;u.u)$, namely $r$.  Apply Tr with identity template and with this
same equation as both premises.  The conclusion is $r=r$.
From a proof of $s=t$, Tr with identity template, major premise $s=s$,
and minor premise $s=t$ gives $t=s$.  Tr then gives transitivity and
term congruence.  These constructions use no introduction rule and no
induction.  In particular, every base singled out by $\mathsf R$ is
already provable in the retained context.

\begin{lemma}[Every case introduction is admissible]
\label{lem:reflexive-all-cases}
For every normal native matrix $B$ and witness $t$,
\[
 \Gamma,B(t,0)\vdash_{\mathcal R}B(t,Sz)
 \quad\Longrightarrow\quad
 \Gamma\vdash_{\mathcal R}D_{x,y}B.
\]
\end{lemma}
\begin{proof}
When $\mathsf R(B,t)$ is false, this is primitive DI.  When it is true,
the discharged base is literally $r=r$.  Use the equality proof above
and ordinary-assumption cut to remove this base from the supplied
derivation.  Weaken the resulting derivation of $B(t,Sz)$ by $B(t,z)$,
then apply the progressive mode of primitive DI.  No substitution or
general progress rule is used.
\end{proof}

\begin{lemma}[One-way preservation is sufficient]
\label{lem:reflexive-substitution}
Arbitrary capture-free normal numerical substitution is admissible in
$\mathcal R$.
\end{lemma}
\begin{proof}
If $B(t,0)$ is literally $r=r$, applying any normal substitution yields
$\operatorname{nf}(r\sigma)=\operatorname{nf}(r\sigma)$.
The compatibility laws identify this with the zero instance of the
substituted matrix and witness.  Thus a true guard remains true.

Now induct on a derivation.  IS, Tr, QE and DE commute with normal
substitution by the normalization facts and the standard fresh-name
convention.  In a progressive DI step the guard was true and remains
true, so the induction hypothesis is the exact premise of progressive
DI at the substituted matrix and witness.  In a case DI step the
induction hypothesis is the substituted case premise.  Apply
Lemma~\ref{lem:reflexive-all-cases}, whether or not its guard has newly
become true.  This completes the induction.
\end{proof}

The reflection direction genuinely fails and is unnecessary.  For
example, an equation with zero instance $p=q$ is initially nonreflexive
when $p,q$ are distinct parameters, but becomes reflexive after
$p:=0,q:=0$.  The substitution proof repairs precisely that change;
it does not falsely claim that every rule label is preserved.

The dependency order so far is normalization laws; weakening and cut;
elementary equality; unrestricted case introduction; arbitrary
substitution.  None of these steps presupposes the equivalence theorem.

\subsection{Recovering progress by equational decoders}
\label{sec:equational-decoders}

The general equation-progress rule can be recovered without constructing
an implication macro and without substituting numerical terms into the
supplied step derivation.  This shorter proof uses only equality,
case-derived closure, weakening, renaming, and ordinary-assumption cut.
All formulas below are zero-normal, and substitutions into formulas mean
the existing normal substitutions.

The needed case-derived closure rule is
\begin{equation}
 \frac{\Lambda,F(0)\vdash F(Sz)}{\Lambda\vdash Q_yF(y)},
 \label{eq:decoder-qcase}
\end{equation}
where $z$ is fresh for $\Lambda$ and the fixed parameters.  It is already
available from unrestricted case DI, followed by DE with an unused
witness.  Unrestricted case DI itself follows from the reflexive-base
introduction: if its guard is true, prove the reflexive base, cut it out,
and weaken by the progressive hypothesis.  Neither construction needs
general equation progress or arbitrary numerical substitution on proofs.

Let
\[
 C(v,a,b):=I(v,a;u.b),\qquad u\notin\operatorname{FV}(b).
\]
Its zero equation is definitional and its successor equation is a direct
instance of the remaining computation rule.

\begin{lemma}[Case decoder]
\label{lem:case-decoder}
Let $r(y)$ be a normal term template and let $a$ be a term independent of
$y$.  Define
\[
 D_{r,a}(y):=
 \bigl(C(y,a,C(y,0,r(y)))=r(y)\bigr).
\]
Then $Q_yD_{r,a}(y)$ is derivable without assumptions.  Consequently
\[
 \Lambda\vdash a=r(0)
 \quad\Longrightarrow\quad
 \Lambda\vdash D_{r,a}(v)
\]
for every term $v$.
\end{lemma}
\begin{proof}
The zero instance of $D_{r,a}$ is literally $a=r(0)$ in zero-normal
syntax.  At $y=Sz$, two successor computations and transitivity prove
\[
 C(Sz,a,C(Sz,0,r(Sz)))
   =C(Sz,0,r(Sz))=r(Sz).
\]
This proof ignores the temporary base hypothesis.  Rule
\eqref{eq:decoder-qcase} therefore gives $Q_yD_{r,a}(y)$ in the empty
context.  Weaken this closure into $\Lambda$ and apply QE with the
given base equation and target $v$.
\end{proof}

This is a valid fixed-context numerical case argument.  The context
$\Lambda$ is held fixed and has no occurrence of the fresh generalized
variable $z$.  The conclusion may later be instantiated at a variable
that appears in other hypotheses because QE, unlike the introduction
rule, imposes no freshness condition on its numerical target.

\begin{theorem}[Unrestricted equation progress]
\label{thm:decoder-progress}
Suppose the supplied derivation is
\[
 \pi:\quad\Gamma,E(z)\vdash E(Sz),\qquad
 E(y)\equiv s(y)=t(y),
\]
with $z$ fresh for $\Gamma$ and the fixed parameters.  Then the
reflexive-base calculus derives $\Gamma\vdash Q_yE(y)$.
\end{theorem}
\begin{proof}
Write $s_0=s(0)$, $t_0=t(0)$, and temporarily retain the original base:
\[
 \Delta:=\Gamma,E(0)=\Gamma,(s_0=t_0).
\]
The progressive lift is
\[
 H(y):=\bigl(C(y,0,s(y))=C(y,0,t(y))\bigr).
\]
Its zero-normal base is literally $0=0$, so it qualifies for primitive
progressive DI.  Define the two decoder equations
\begin{align*}
 D_s(y)&:\quad C(y,s_0,C(y,0,s(y)))=s(y),\\
 D_t(y)&:\quad C(y,s_0,C(y,0,t(y)))=t(y).
\end{align*}
Lemma~\ref{lem:case-decoder} gives their closures without assumptions.
For $D_s$ the zero base is $s_0=s_0$, supplied by reflexivity; for
$D_t$ it is $s_0=t_0$, supplied by the retained hypothesis $E(0)$.
Hence, in the fixed context $\Delta$, QE gives both $D_s(z)$ and
$D_t(z)$ at the fresh variable $z$.

Now add $H(z)$.  Congruence in the one-term context
\[
 w\longmapsto C(z,s_0,w)
\]
converts it to
\[
 C(z,s_0,C(z,0,s(z)))
   =C(z,s_0,C(z,0,t(z))).
\]
Symmetry of $D_s(z)$ followed by this equation and $D_t(z)$ yields
$s(z)=t(z)$, that is, $E(z)$.

Weaken the original step derivation $\pi$ into this enlarged retained
context, and cut the just-derived $E(z)$ into its distinguished
hypothesis.  The result is $\Delta,H(z)\vdash E(Sz)$.
No substitution into $\pi$ has occurred: its induction variable is
still the same $z$.  Successor computation on each side of $H(Sz)$
then converts this to
\[
 \Delta,H(z)\vdash H(Sz).
\]
Use an unused witness coordinate and witness $0$ in primitive
progressive DI.  The conclusion is $\Delta\vdash D_{x,y}H(y)$;
DE returns $\Delta\vdash Q_yH(y)$ because the witness is unused.
Reflexivity proves $H(0)=(0=0)$, and QE at $Sz$ gives $H(Sz)$, hence
$E(Sz)$.  We have therefore constructed
\[
 \Gamma,E(0)\vdash E(Sz).
\]
Finally rule~\eqref{eq:decoder-qcase} discharges the retained base and
gives $\Gamma\vdash Q_yE(y)$.
\end{proof}

The complete dependency order is especially short: normal-substitution
identities; weakening and cut; elementary equality; unrestricted case
DI; case-derived closure; the two decoders; general equation progress.
Arbitrary numerical substitution on derivations can consequently be
proved afterward, or transferred through the exact comparison with the
seven-rule zero-normal calculus.  The decoder proof does not use that
structural theorem as an unproved prerequisite.

\paragraph{Relation to the executable translation.}
In \texttt{proof\_translations.py}, \texttt{decoder\_closure} constructs
the two case-derived closure proofs.  In \texttt{qi\_to\_six}, their QE
instances are \texttt{dl} and \texttt{dr}; \texttt{decoded} is the
congruence-and-transitivity derivation of $E(z)$; and \texttt{cut\_tree}
inserts it into the supplied step proof.  The final two operations are
guarded progress for $H$, then case-derived closure for $E$.
Proof-tree numerical lifting merely renames free parameters away from
new eigenvariables.  It is not arbitrary term substitution into the
input step derivation.


\subsection{Exact equivalence and bi-interpretation}

\begin{theorem}[Same normal sequents]
For every context $\Gamma$ and formula $A$ of the zero-normal language,
\[
 \Gamma\vdash_{\mathcal R}A
 \quad\Longleftrightarrow\quad
 \Gamma\vdash_{\mathcal Z}A.
\]
Both transformations are effective on finite derivations and are the
identity on the sequent syntax.
\end{theorem}
\begin{proof}
For $\mathcal R\to\mathcal Z$, the first five rules are shared.
A case-selected DI is native case introduction.  If DI is
progress-selected, its matrix must be an equation.  Apply the native
equation-progress rule to $E(y):=B(t,y)$ to obtain $Q_yE(y)$.
Under the temporary assumption $E(0)$, QE supplies $E(Sz)$.
The native D case rule discharges $E(0)$ and produces the required
$D_{x,y}B$.  Hence every primitive rule is sound in $\mathcal Z$.

For $\mathcal Z\to\mathcal R$, translate recursively, retain the
five shared rules, replace case introduction by
Lemma~\ref{lem:reflexive-all-cases}, and replace equation progress by
Theorem~\ref{thm:decoder-progress}.  Each construction preserves the
retained context and its freshness conditions.
\end{proof}

Compose this exact normal-sequent equivalence with the repository's
$\mathcal Z$--native normalization comparison and the established
native--iPRA bi-interpretation.  This yields a bi-interpretation and
equiconsistency with full first-order iPRA.  Both numerical domains
remain whole and interpreted equality remains actual equality; the
internal numerical round-trip graph is $x=y$.  Formula round trips are
the inherited round trips composed with the provable equivalence of a
formula and its zero-normal form.  The normalization layer is explicitly
part of the presentation, and is included in any checker measurement.

\subsection{The same guard on raw syntax}

\begin{proposition}[Seven rules without normalization]
On the original raw iterator syntax, retain Hyp, I0, IS, Tr, QE and DE,
and replace the two introductions by $\mathrm{DI}_{\mathsf R}$, with
$\mathsf R(B,t)$ now testing the unnormalized literal syntax of
$B(t,0)$.  The resulting seven-rule calculus derives exactly the
sequents of the original eight-rule native kernel.
\end{proposition}
\begin{proof}
Use ordinary capture-free raw substitution throughout.  Reflexivity
comes from I0 and crossed Tr: apply identity-template transport to the
equation $I(0,r;u.u)=r$ twice.  Weakening and cut are direct inductions,
every missing case introduction is recovered by reflexivity and cut,
and true guards are preserved by substitution.  The false-to-true repair
therefore proves arbitrary substitution exactly as above.  The case-derived closure rule is recovered exactly as above.

The one modification is the lift used to recover unrestricted progress.
For $E(y)=(s(y)=t(y))$, define the asymmetric equation
\[
 K(y):=
 \bigl(C(y,t(y),s(y))=C(0,t(y),s(y))\bigr),
 \qquad C(v,a,b):=I(v,a;u.b),
\]
with each accumulator binder fresh.  Its raw zero instance is literally
\[
 C(0,t(0),s(0))=C(0,t(0),s(0)),
\]
so it satisfies the reflexive-base guard before any computation.
For every $v$, IS computes the left side of $K(Sv)$ to $s(Sv)$,
while I0 computes its right side to $t(Sv)$.  Elementary equality
therefore gives $K(Sv)\dashv\vdash E(Sv)$.

Under the retained base $s(0)=t(0)$, the raw decoder equations are
\begin{align*}
 C(y,s(0),C(y,t(y),s(y)))&=s(y),\\
 C(y,s(0),C(0,t(y),s(y)))&=t(y).
\end{align*}
Their zero cases follow from I0 and respectively reflexivity or the retained
base.  Their successor cases follow from IS for the outer selector and
IS or I0 for the inner selector.  Case-derived closure and QE supply both
decoders.  Congruence of $K(y)$ under $C(y,s(0),-)$ therefore yields $E(y)$.
Cut this into the unchanged supplied progress proof, convert $E(Sy)$
to $K(Sy)$, and finish with primitive progress followed by case-derived
closure, as in Theorem~\ref{thm:decoder-progress}.  The proof translations
then give the same raw sequents without a normalization map.
\end{proof}

The asymmetric lift also works in the zero-normal presentation, where
its base contracts to $t(0)=t(0)$.  Thus the broad reflexive-base guard
is independent of the choice to make zero computation definitional.
This raw proof does not establish the narrower exact-$0=0$ variant on
raw syntax: the displayed raw base is reflexive but is generally not
literally $0=0$.

\subsection{Further variants and the implementation tradeoff}

\paragraph{Exact-zero base.}
One may replace $\mathsf R(B,t)$ by the smaller class
\[
 \mathsf Z(B,t)\quad\Longleftrightarrow\quad
 B(t,0)\text{ is literally }0=0.
\]
Every proof above remains valid: selected bases are provable, a true
guard is preserved by substitution, and every gate $H$ in the progress
proof has exactly this base.  Thus this variant also derives precisely
the native zero-normal sequents.  It uses a single comparison with a
fixed formula, though it recognizes fewer usable progress matrices.

\paragraph{Testing before the witness is inserted.}
Alternatively test whether $B(x,0)$, with $x$ still a formal variable,
is reflexive.  The same argument works: a true guard makes every
$B(t,0)$ reflexive, arbitrary gate matrices still qualify, and true
guards remain true under substitution of the retained parameters.
This removes the witness from the test but usually costs an extra
normal-substitution pass in a checker, because the instantiated base
$B(t,0)$ must still be constructed for the unguarded local assumption.
The post-witness test above reuses that already-needed base.

\paragraph{Scope of the contribution.}
The guarded-introduction completeness method is inherited from the
repository's diagonal-guard package.  The new specialization tests a
computed zero instance rather than an exact iterator shape or a literal
witness, and its structural proof explicitly handles a guard that can
change from false to true.  It is not a new arithmetic bridge or a
proof of global minimality.  A reduction of mathematical bookkeeping
does not by itself imply a new record in all checker-size coordinates.

\paragraph{A distinct inherited rule must not be claimed as new.}
Retaining both $B(t,0)$ and, at equations, $B(t,z)$ in a single
introduction premise is another short equivalent presentation.  That
cumulative introduction already appears as the old DI in the
inherited Choice-DI material.  It may be used as a control or composed
with zero normalization, but it is not a new result of this section.
